\scanpage{6}
\mspage{4}

dans $\mathbb P^3_k$, alors on a
\[
 \mathbb Z\simeq\operatorname{Pic}(\mathbb P^3)\xrightarrow{\sim}\operatorname{Pic}(X_0),
\]
$k$ étant un corps de car. quelconque.

2. \uline{Rappels sur la structure de} $\pi$.

On suppose $S$ hensélien, i.e. $V$ hensélien.

2.1. Il y a une sous-extension $\widetilde K_{\deleted{\unsure{u}}}$ de $\bar K/K$, qui est la plus grde ss-extension non ramifiée (i.e. étale) \deleted{\unsure{dans}} sur $V$. Si \deleted{$\unsure{\widetilde K_\ell}$ est} $\bar V$ est la clôture normale de $V$ dans $\bar K$, $\widetilde V=\deleted{\unsure{\bar V_\ell}}\bar V\cap\widetilde K_{\deleted{\unsure{u}}}$ la clôture normale dans $\widetilde K_{\deleted{\unsure{u}}}$, alors $\widetilde V_{\deleted{\unsure{u}}}$ est un anneau de valuation discrète tel que
\[
 \deleted{\unsure{\mathfrak m}\widetilde V=\mathfrak m\widetilde V,}
 \qquad \widetilde V\otimes_V k=\bar k
\]
est une \deleted{extension de} clôture séparable de $k$; En d'autres termes, $\widetilde V$ est une \uline{hensélisation stricte} de $V$ p. r. à \deleted{\unsure{une}} la \deleted{\unsure{sa}} clôture séparable $\bar k$ de $k$, et $\widetilde K_{\deleted{\unsure{u}}}$ est le corps des fractions de $\widetilde V$. On pose
\inserted{
\[
\begin{cases}
\widetilde S=\operatorname{Spec}\widetilde V=\text{hensélisé strict de }S\\
\unsure{\bar s}=\operatorname{Spec}\bar k.
\end{cases}
\]
}

\margin{\figure{\begin{tikzpicture}[x=1cm,y=1cm]
\draw (0,0) circle(0.06);\draw (0.06,0)--(0.85,0.02);
\draw (1.02,0.85)..controls(0.9,0.35) and(1.03,0.2)..(1.4,0.08)..controls(1.07,-0.1) and(1.02,-0.35)..(0.94,-0.8);
\end{tikzpicture}}}
\unsure{2.3|2.2}. On se rappelle que \deleted{\unsure{toute}} la donnée de la clôture séparable $\bar K$ de $K$ définit canoniquement une clôture séparable $\bar k$ de $k$.

\deleted{Pour \deleted{\unsure{tout}} \unsure{$n$} \deleted{\unsure{premier à}} \unsure{$p$} \inserted{Si \unsure{$n$} est un} entier \unsure{$n$ premier à} \unsure{$p$} inversible dans $k$, on a
\[
2.1.1.\qquad {}_n\bar k^*\xleftarrow[\unsure{\varphi}]{\sim}{}_n\widetilde V^*\xrightarrow{\sim}\bar K^*
\]
d'où un isomorphisme canonique
\[
\mu_n(\bar K)\simeq\mu_n(\bar k).
\]
Ces isom. sont compatibles avec les \unsure{empilements}
}
\margin{5}
